\input{../../../templates/es/preambulo.tex}

% Los listados se toman de src/ con \lstinputlisting; la diapositiva es el archivo.
\lstset{
    basicstyle=\ttfamily\fontsize{8pt}{9.6pt}\selectfont,
    breaklines=true,
    breakatwhitespace=true,
    columns=fullflexible,
    keepspaces=true,
    xleftmargin=0.3em,
    framesep=2pt,
    aboveskip=0pt,
    belowskip=0pt
}

\newcommand{\marcoresultado}[3]{%
    \begin{frame}{#1}
        \begin{center}
            \includegraphics[height=0.62\textheight,width=\textwidth,keepaspectratio]{#2}
        \end{center}
        \vspace{0.25em}
        {\small #3\par}
    \end{frame}%
}

\newcommand{\marcoafirmacion}[2]{%
    \begin{frame}{#1}
        {\small #2\par}
    \end{frame}%
}

\date{}

\begin{document}

\tituloleccion{05}{Mutación}

\section{Esta lección}

\begin{frame}{Esta lección}
La Lección 04 partió la cruza en dos familias que no se pueden intercambiar. La mutación es el operador que queda, todavía medida sobre un cromosoma.

\vspace{0.6em}
\begin{itemize}
    \item El alcance es una tasa por un paso; hay que ver las dos.
    \item Un cromosoma de permutación no admite ruido gaussiano.
    \item Un óptimo sin volumen es un artefacto que ninguna búsqueda puede ocupar.
\end{itemize}
\end{frame}

% ================= alcance_mutacion_01_el_instrumento =================
\begin{frame}{Un núcleo de mutación}
Para el gen \(i\),
\[
B_i\sim\operatorname{Bernoulli}(p),\qquad
\Delta_i\sim\mathcal N(\mu,\sigma^2),\qquad
x_i'=\Pi_{[l,u]}(x_i+B_i\Delta_i).
\]
Entonces \(\mathbb E[\sum_iB_i]=Lp\): la tasa controla cuántos genes se
activan y \(\sigma\) el desplazamiento. Se miden posiciones cambiadas con
\(H=\sum_i\mathbf 1[x_i\ne x_i']\) y recorrido de permutación con
\(D=\sum_v|\operatorname{pos}_\pi(v)-\operatorname{pos}_{\pi'}(v)|\). Un óptimo de ancho cero tiene
probabilidad cero bajo una propuesta continua ideal, aunque una cuadrícula lo
contenga.
\end{frame}

\section{El alcance de una mutación, y cómo medirlo}

\begin{frame}{El alcance de una mutación, y cómo medirlo}
La selección descarta y la cruza recombina; ambos trabajan solo con valores que la
población ya posee. La mutación es el único operador que puede producir un valor
que nadie tenía. Por lo tanto, toda esta lección trata sobre una pregunta: ¿QUÉ TAN LEJOS
puede llegar una mutación?
\end{frame}

\marcoresultado{El alcance de una mutación, y cómo medirlo}{../figures/alcance_mutacion_01_el_instrumento.png}{%
    Los dos números de la Lección 01 reaparecen exactamente — sigma 1.0 llega \textbf{0 de 2000}, sigma 3.0 \textbf{110 de 2000}. El viaje promedio es 0.798·sigma hasta que los muros empiezan a comerlo: en sigma 6.0, el \textbf{19.4\%} de las propuestas son acotadas y la proporción baja a 0.690. El intercambio en una tabla: sigma 0.5 encuentra la mayoría de los movimientos cuesta arriba (\textbf{272}) y llega \textbf{0} veces}

% ================= alcance_mutacion_02_tasa_y_paso =================
\section{La mutación tiene dos diales, no uno}

\begin{frame}{La mutación tiene dos diales, no uno}
El paso 1 midió un solo gen, por lo que sigma era lo único que había para ajustar.
Un cromosoma real tiene muchos genes, y el operador tiene entonces un segundo dial: la
probabilidad por gen p de que un gen sea tocado en absoluto. El
\texttt{mutation\_random\_deviation(ind, mu, sigma, p)} de Gridin lleva ambos.
\end{frame}

\begin{frame}[fragile]{(1) CANTIDAD\_GENES}
\lstinputlisting[basicstyle=\ttfamily\fontsize{8pt}{9.6pt}\selectfont,firstline=37,lastline=43]{../src/alcance_mutacion_02_tasa_y_paso.py}
\end{frame}

\begin{frame}[fragile]{(2) Individuo.aptitud}
\lstinputlisting[basicstyle=\ttfamily\fontsize{8pt}{9.6pt}\selectfont,firstline=65,lastline=65]{../src/alcance_mutacion_02_tasa_y_paso.py}
\end{frame}

\begin{frame}[fragile]{(3) mutar\_desviacion\_aleatoria()}
\lstinputlisting[basicstyle=\ttfamily\fontsize{8pt}{9.6pt}\selectfont,firstline=81,lastline=96]{../src/alcance_mutacion_02_tasa_y_paso.py}
\end{frame}

\begin{frame}[fragile]{(4) reporte\_regimen()}
\lstinputlisting[basicstyle=\ttfamily\fontsize{6pt}{7.2pt}\selectfont,firstline=99,lastline=135]{../src/alcance_mutacion_02_tasa_y_paso.py}
\end{frame}

\marcoresultado{La mutación tiene dos diales, no uno}{../figures/alcance_mutacion_02_tasa_y_paso.png}{%
    Los genes tocados nunca difieren de \texttt{p·n} por más de \textbf{0.015}. Tres regímenes con el mismo \texttt{p·sigma = 0.30} viajan 1.904 / 2.179 / 2.304 en total pero su salto individual más largo difiere por un factor de \textbf{7.1} (10.253 vs 1.449). En p = 0.1, \textbf{718 de 2000} mutaciones no cambian nada en absoluto, contra un predicho (1−p)¹⁰ = 0.3487}

% ================= alcance_mutacion_03_guiada_por_aptitud =================
\section{Mutación que se niega a empeorar las cosas}

\begin{frame}{Mutación que se niega a empeorar las cosas}
Cada mutación hasta ahora ha sido ciega: mueve los genes y acepta lo que sea que
diga el objetivo sobre el resultado. La mayor parte del tiempo el resultado es peor -
el paso 1 ya mostró que en sigma = 3.0 solo 144 de 2000 mutaciones fueron
cuesta arriba. El último operador de Gridin las rechaza: reintenta, hasta MAX\_INTENTOS
veces, y devuelve el primer mutante que supera a su padre.
\end{frame}

\begin{frame}[fragile]{(1) MAX\_INTENTOS}
\lstinputlisting[basicstyle=\ttfamily\fontsize{8pt}{9.6pt}\selectfont,firstline=39,lastline=44]{../src/alcance_mutacion_03_guiada_por_aptitud.py}
\end{frame}

\begin{frame}[fragile]{(2) mutar\_guiada\_por\_aptitud()}
\lstinputlisting[basicstyle=\ttfamily\fontsize{7pt}{8.5pt}\selectfont,firstline=95,lastline=119]{../src/alcance_mutacion_03_guiada_por_aptitud.py}
\end{frame}

\begin{frame}[fragile]{(3) comparar()}
\lstinputlisting[basicstyle=\ttfamily\fontsize{6pt}{7.2pt}\selectfont,firstline=122,lastline=165]{../src/alcance_mutacion_03_guiada_por_aptitud.py}
\end{frame}

\marcoresultado{Mutación que se niega a empeorar las cosas}{../figures/alcance_mutacion_03_guiada_por_aptitud.png}{%
    El filtro convierte un cambio de aptitud promedio de \textbf{−1.0337 en +0.0548} y duplica las llegadas (102 → 273) — pero gasta \textbf{5,603 evaluaciones de aptitud contra 2,000}, así que por mil llamadas llega \textbf{48.7} veces contra \textbf{51.0} de la mutación ciega. La aceptación del 20.2\% es exactamente 1−(1−q)³ para la tasa ciega hacia arriba q}

% ================= alcance_mutacion_04_la_escalera_muerta =================
\section{La ejecución que la Lección 02 dio por muerta}

\begin{frame}{La ejecución que la Lección 02 dio por muerta}
Una versión anterior del paso 5 de la Lección 02 terminó con una ejecución que se detuvo a las seis
generaciones, a 0.2000 por debajo del óptimo de la escalera, casi sin dispersión genética
restante, y apuntó hacia aquí: la mutación "es el único operador que podría restaurarlo". Este paso es donde esa promesa
venció - y no venció de la forma
en que estaba escrita. Lo que la investigación encontró en su lugar envió una corrección de vuelta a
la Lección 02, que ya no hace esa afirmación.
\end{frame}

\begin{frame}[fragile]{(1) ejecutar()}
\lstinputlisting[basicstyle=\ttfamily\fontsize{6pt}{7.2pt}\selectfont,firstline=115,lastline=201]{../src/alcance_mutacion_04_la_escalera_muerta.py}
\end{frame}

\begin{frame}[fragile]{(2) escalera()}
\lstinputlisting[basicstyle=\ttfamily\fontsize{8pt}{9.6pt}\selectfont,firstline=93,lastline=107]{../src/alcance_mutacion_04_la_escalera_muerta.py}
\end{frame}

\begin{frame}[fragile]{(3) anchos\_escalones()}
\lstinputlisting[basicstyle=\ttfamily\fontsize{7pt}{8.5pt}\selectfont,firstline=204,lastline=223]{../src/alcance_mutacion_04_la_escalera_muerta.py}
\end{frame}

\begin{frame}[fragile]{(4) ALTURA\_PICO\_REPARADA}
\lstinputlisting[basicstyle=\ttfamily\fontsize{8pt}{9.6pt}\selectfont,firstline=226,lastline=233]{../src/alcance_mutacion_04_la_escalera_muerta.py}
\end{frame}

\begin{frame}[fragile]{(5) barrido()}
\lstinputlisting[basicstyle=\ttfamily\fontsize{6pt}{7.2pt}\selectfont,firstline=236,lastline=264]{../src/alcance_mutacion_04_la_escalera_muerta.py}
\end{frame}

\marcoresultado{La ejecución que la Lección 02 dio por muerta}{../figures/alcance_mutacion_04_la_escalera_muerta.png}{%
    La ejecución de la Lección 02 reproducida (6 generaciones, mejor \textbf{+1.8000}, dispersión 0.192, \textbf{0.2000} por debajo). Luego ocho regímenes de mutación × 100 ejecuciones: el escalón superior se alcanza \textbf{0 veces en 800 ejecuciones}, mientras que la dispersión promedio de los genes va de 0.462 a 3.194. \texttt{anchos\_escalones()} dice por qué — el escalón superior es \textbf{0.000 de ancho}. Reparada (altura del pico 10.0 → 10.5, escalón superior 0.500 de ancho), el propio régimen de la Lección 02 lo alcanza \textbf{96 de 100} veces}

% ================= mutacion_permutacion_01_ruido_ilegal =================
\section{cuando el ruido no es un movimiento legal}

\begin{frame}{cuando el ruido no es un movimiento legal}
Cada operador hasta ahora ha sumado un número a un gen. Eso funciona porque un gen
era un valor real y cualquier valor real en [-10, +10] era un gen legal.
\end{frame}

\marcoresultado{cuando el ruido no es un movimiento legal}{../figures/mutacion_permutacion_01_ruido_ilegal.png}{%
    La desviación aleatoria en una permutación produce \textbf{1935 mutantes inválidos de 2000}, y \textbf{0} de los 65 sobrevivientes tuvo algún gen tocado — la única salida legal del operador es la intacta, a una tasa medida que coincide con (1−p)¹⁰ a menos de \textbf{0.0025}. La inversión de bit cambia exactamente 1 posición cada vez; el intercambio exactamente 2}

% ================= mutacion_permutacion_02_reordenamientos =================
\section{tres reordenamientos más grandes}

\begin{frame}{tres reordenamientos más grandes}
El intercambio es el movimiento legal mínimo. Gridin da tres más grandes, y los tres
funcionan de la misma manera: elige dos posiciones y reordena el bloque entre
ellas. La inversión lo invierte (le da la vuelta), el desplazamiento lo rota por un lugar, barajar
lo aleatoriza. La pregunta que responde este paso es la que el primer ejemplo hizo sobre
sigma - ¿qué tan lejos llega cada uno de ellos? - y la respuesta necesita dos números,
no uno.
\end{frame}

\begin{frame}[fragile]{(1) mutacion\_inversion()}
\lstinputlisting[basicstyle=\ttfamily\fontsize{7pt}{8.5pt}\selectfont,firstline=79,lastline=97]{../src/mutacion_permutacion_02_reordenamientos.py}
\end{frame}

\begin{frame}[fragile]{(2) mutacion\_desplazamiento()}
\lstinputlisting[basicstyle=\ttfamily\fontsize{7pt}{8.5pt}\selectfont,firstline=100,lastline=119]{../src/mutacion_permutacion_02_reordenamientos.py}
\end{frame}

\begin{frame}[fragile]{(3) mutacion\_barajar()}
\lstinputlisting[basicstyle=\ttfamily\fontsize{7pt}{8.5pt}\selectfont,firstline=122,lastline=140]{../src/mutacion_permutacion_02_reordenamientos.py}
\end{frame}

\begin{frame}[fragile]{(4) comparar\_operadores()}
\lstinputlisting[basicstyle=\ttfamily\fontsize{6pt}{7.2pt}\selectfont,firstline=167,lastline=199]{../src/mutacion_permutacion_02_reordenamientos.py}
\end{frame}

\marcoresultado{tres reordenamientos más grandes}{../figures/mutacion_permutacion_02_reordenamientos.png}{%
    Las dos medidas de alcance clasifican a los operadores de manera diferente: por posiciones perturbadas \textbf{desplazamiento (4.663) > inversión (4.230) > barajar (3.708) > intercambio (2.000)}, por distancia viajada \textbf{inversion (1.308) > barajar (0.876) > intercambio (0.733) = desplazamiento (0.733)}. El intercambio y el desplazamiento dan idéntico viaje total en \textbf{90 de 90} pares de posiciones. Barajar devuelve el individuo sin cambios \textbf{271 de 2000} veces (13.6\%)}

\section{Siguiente}

\begin{frame}{Siguiente}
La escalera de la Lección 02 tuvo un peldaño de arriba de ancho cero. Esta lección conserva la geometría rota como el ejemplar que diagnostica.

\vspace{0.8em}
\textbf{Lección 06 --- Medir la efectividad}

Un algoritmo genético es una variable aleatoria. Una generación, y una corrida, no prueban nada. La Lección 03 aplazó esa evidencia aquí.
\end{frame}
\end{document}
